\documentclass[conference]{IEEEtran}
\IEEEoverridecommandlockouts
% The preceding line is only needed to identify funding in the first footnote. If that is unneeded, please comment it out.
\usepackage{cite}
\usepackage{amsmath,amssymb,amsfonts}
\usepackage{algorithmic}
\usepackage{graphicx}
\usepackage{textcomp}
\usepackage{xcolor}
\usepackage{url}
\def\BibTeX{{\rm B\kern-.05em{\sc i\kern-.025em b}\kern-.08em
    T\kern-.1667em\lower.7ex\hbox{E}\kern-.125emX}}
\begin{document}

\title{Replication and Automated Evaluation of QUIS: Studying Structured Question Generation for Exploratory Data Analysis}

% Author information removed for double-blind review.
\author{\IEEEauthorblockN{Anonymous Authors}
\IEEEauthorblockA{Anonymous Institution\thanks{Code and evaluation scripts: \protect\url{https://github.com/sunnydovision/EDAProject}}}
}

\maketitle

\begin{abstract}
This paper is a replication study of QUIS, an automated Exploratory Data Analysis (EDA) system. QUIS uses a question generation module (QUGEN) to create questions that are later used as a guideline for searching for insights. The original QUIS paper used human evaluation, but this is difficult to reproduce and also takes much time. In this work, we propose an automatic evaluation framework with 9 metrics covering correctness, structural validity, statistical quality, subspace exploration, and question-insight alignment. We compare our QUIS replication with two baselines: an ablation called ONLYSTATS and an LLM-based agentic baseline that we implemented. We evaluate three systems across three datasets: Adidas US Sales, IBM Employee Attrition, and Online Sales. The experiment results show that QUGEN helps avoid column-type errors and improves subspace exploration. Across all datasets, QUIS obtains 94.0\% average structural validity, while the LLM baseline obtains only 40.0\%. QUIS also produces more insights with subspace filters: 84.4\% on average, compared with 37.4\% for the LLM baseline. For subspace quality, QUIS achieves the highest Score Uplift ($x$=1.067) and Simpson's Paradox Rate (30.1\%). We observe that structured question generation improves both exploration breadth and the quality of discovered subspace patterns.
\end{abstract}

\begin{IEEEkeywords}
Exploratory Data Analysis, Automated EDA, Question Generation, Evaluation Metrics, Comparative Evaluation
\end{IEEEkeywords}

\section{Introduction}
Basically, automated EDA requires expert's knowledge and time to explore hidden patterns and relationships in new datasets. As datasets grow, automated EDA systems become more valuable. QUIS (Manatkar et al., 2024) \cite{quis} is one such system: rather than searching for insights directly, it generates analysis questions through a module called QUGEN, then uses a second module named ISGEN to search the dataset for matched insights corresponding to each question.

For the original QUIS paper, evaluation was done through human ratings, which are costly, subjective, and difficult to reproduce. Besides, since the original QUIS implementation is not published, we replicate QUIS based on the system description in the original paper. Our goal is not to verify numerical reproducibility, but to use this replication as an instrument for studying the contribution of each pipeline component through automated metrics. We also design an evaluation study to answer two questions:

\begin{enumerate}
\item Can we build an automatic evaluation framework for comparing automated EDA systems?
\item What does the QUIS pipeline contribute---specifically, QUGEN (card generation) and ISGEN (insight search)---when compared with simpler baselines?
\end{enumerate}

We replicate QUIS and compare it with two baselines across three datasets: Adidas US Sales, IBM Employee Attrition, and Online Sales. The first baseline is ONLYSTATS, an ablation of QUIS that removes QUGEN and uses a statistics-based selection method for card generation. The second one is an LLM-based agentic EDA pipeline. We use a total of 9 metrics to study correctness, structure, statistical quality, subspace exploration, and question quality.

Our contributions are:

\begin{itemize}
\item We build an automatic and human-free evaluation framework with 9 metrics for AutoEDA systems, covering correctness, structural validity, statistical quality, and subspace exploration.
\item We indicate that QUGEN's card generation prevents systematic column type mismatches that trouble LLM pipelines (SVR: Baseline's 40.0\% vs QUIS's 94.0\%).
\item We prove that ISGEN's beam search, guided by QUGEN's insight cards, achieves better subspace coverage (QUIS's 84.4\% vs Baselines' 37.4\%) and quality (Score Uplift $x$=1.067, SPR 30.1\%).
\item We show that structural validity is a prerequisite for meaningful subspace discovery: the LLM Baseline, with only 40.0\% SVR, produces zero significant paradoxes across all datasets, while QUIS produces significant paradoxes in two out of three.
\end{itemize}

\section{Related Work}

\subsection{Automated EDA and Insight Evaluation}

Early automated EDA systems measure fascination using statistical deviation. SEEDB \cite{seedb} recommends visualizations by comparing distributions between a target subset and a reference, ranking candidates by the deviation score. QuickInsights \cite{quickinsights} develops a unified insight model covering multiple pattern types: trend, outlier, attribution, then scores each by combined impact and significance for efficient mining. Zenvisage \cite{zenvisage} takes a query-based approach, allowing users to specify desired visual patterns and retrieve matching views from whole collections. These systems enable automatic ranking without human input, but their quality metrics are limited and do not validate whether an insight has the correct structural form (for example: a trend computed over a categorical column instead of a temporal one still produces a non-zero deviation score).

Visualization-focused systems assess quality by how well they look and how clearly they organize information. Voyager \cite{voyager} automatically generates and ranks views using expressiveness and effectiveness rules that are derived from perceptual principles. Voder \cite{voder} combines natural language facts with interactive visuals to highlight interesting findings during data exploration. MetaInsight \cite{metainsight} organizes insights into structured knowledge by measuring the precise patterns that can be extracted from data. These approaches assess visualization design quality and exploration coverage, but do not check whether the breakdown column semantics match the intended pattern type.

More recent systems, specifically, InsightPilot \cite{insightpilot}, uses LLMs to suggest analysis actions and evaluate completeness through user studies. QUIS \cite{quis} generates questions from the dataset before searching for insights, and evaluates output quality using human rater scores on a 1--5 scale. Not only human evaluation captures subjective quality, it also has limitations: expensive, time-consuming, and difficult to reproduce. These limitations motivate the need for automatic evaluation frameworks that can be applied consistently across systems.

\subsection{Subgroup Analysis and Pattern Quality}

Beyond global statistics, recent work proliferously focuses on conditional insights---patterns found within specific data subsets. Subgroup discovery \cite{bach} formalizes the finding of data subsets with notable target behavior, using quality measures such as weighted relative accuracy (WRAcc) to calculate how much a subgroup deviates from the overall population. A key question in this line of work is whether a discovered subgroup reveals stronger or weaker patterns than the global baseline---Score Uplift measuring how much stronger patterns are in a subspace compared to globally.

A related and outstanding statistical study is Simpson's Paradox \cite{simpson}, where an association that holds overall flips direction when the data is broken into subgroups. Simpson's original work showed that merging contingency tables can produce conclusions that contradict those from each subgroup individually. Teng and Lin \cite{teng} examine Simpson's Paradox in observational data, proposing a kernel-based tree algorithm to detect and explain subgroup reversals driven by confounding and effect heterogeneity. Their work formalizes when subgroup associations diverge from global trends, framing such reversals as analytically meaningful rather than statistical noise. However, no prior work has combined subgroup significance testing and paradox detection into a unified evaluation framework for automated EDA systems.

\subsection{Research Gap}

No existing AutoEDA evaluation framework checks structural validity---whether the breakdown column type matches the intended pattern---or combines subgroup significance testing with paradox detection. \textbf{Our work} addresses both gaps with a fully automatic, human-free framework of 9 metrics covering correctness, structural validity, statistical quality, subspace exploration, and question quality. Our evaluation shows that LLM baselines have high systematic structural errors and produce zero significant paradoxes, while structured approaches achieve the highest SVR and discover meaningful pattern reversals.

\section{System Overview}

We compare three systems in this study.

\subsection{QUIS (Our Replication)}

QUIS has two main stages. First, the QUGEN module reads the dataset schema and basic statistics, then creates Insight Cards using an LLM. Each card contains 4 tuples: a question, a reason, a breakdown column (B), and a measure (M). Second, ISGEN uses these cards and searches the dataset for actual insights.

Each insight has a 4-tuple \textbf{(B, M, S, P)}, where B is the breakdown column, M is the measure, S is the subspace filter (a set of column-value equality conditions), and P is the pattern type. As the origin, we study four pattern types: TREND, OUTSTANDING\_VALUE, ATTRIBUTION, and DISTRIBUTION\_DIFFERENCE.

For each insight card, ISGEN computes the global view ($S = \emptyset$), then runs a beam search to find subspaces where the pattern is stronger. Filter columns in beam search can be selected randomly or via LLM suggestion (probability 0.5). It scores each candidate using pattern-specific functions (Mann-Kendall $\tau$ for TREND, attribution share for ATTRIBUTION, Jensen-Shannon divergence for DISTRIBUTION\_DIFFERENCE), and keeps the top candidates. This subspace search is shared between QUIS and ONLYSTATS---the only difference is the source of the input cards.

\subsection{ONLYSTATS (Ablation Baseline)}

ONLYSTATS is the ablation baseline derived from QUIS. It replaces QUGEN with a statistics-based selection method: it computes Kruskal-Wallis test scores \cite{kruskal} for all possible breakdown-measure (B,M) pairs and then selects the top 20 pairs per breakdown. It produces 140 (B, M) pairs on Adidas, 320 on Employee Attrition, and 60 on Online Sales. Insight Cards are generated from these (B,M) pairs using a template-based approach without LLM. ONLYSTATS then runs the ISGEN pipeline, same as in QUIS. This ablation isolates the contribution of QUGEN to card quality alone.

\subsection{LLM Agentic Baseline}

The second baseline is an LLM-based agentic EDA pipeline that we implemented using \texttt{gpt-5.4} at temperature 0.3--0.5. It consists of five agents: (1) data profiling, which classifies each semantic column; (2) quality analysis, which detects missing values, outliers, and duplicates; (3) statistical analysis, which computes statistics and Pearson correlations; (4) pattern discovery, which searches temporal, correlation, grouping, and anomaly patterns from pre-computed aggregations; and (5) insight extraction, which assembles all prior outputs into structured insights. Each agent uses outputs from previous steps as context. This pipeline does not enforce a structured Insight Card format; the LLM reasons freely over pre-computed evidence. For evaluation, outputs are forced to be processed into a QUIS-compatible format (breakdown column, measure, subspace, pattern type).

\section{Evaluation Framework}

\subsection{Selected Metrics}

We selected 9 essential metrics from a larger list of candidates. We group them as follows.

\subsubsection{Group 1: Correctness}

\textit{Faithfulness} checks whether the numbers in an insight can be recomputed exactly from the dataset. For example, an insight reports a total sales value for one region, we recompute the aggregation from the CSV file to check the precision of the insight report. This metric works as a sanity check: a system with low faithfulness is hallucinating, making all other metrics unreliable. A high score is expected for good systems and confirms that the following metrics measure real patterns.

The formula is:
\begin{equation}
\text{Faithfulness} = \frac{|\{i \in I : \text{verified}(i)\}|}{|I|}
\end{equation}

where $I$ is the set of all insights, and $\text{verified}(i)$ returns true if for all reported values $v_r$ in insight $i$, the recomputed value $v_c$ satisfies $|v_r - v_c| < \epsilon$ with $\epsilon = 10^{-6}$.

\subsubsection{Group 2: Structural Quality}

\textit{Structural Validity Rate (SVR)} checks whether the breakdown column type matches the pattern type. For example, a TREND insight should use a date or time column as the breakdown. An ATTRIBUTION insight should use a categorical breakdown, not a numeric measure column.

The formula is:
\begin{equation}
\text{SVR} = \frac{|\{i \in I : \text{valid\_breakdown}(i)\}|}{|I|}
\end{equation}

where $\text{valid\_breakdown}(i)$ checks pattern-specific rules: for TREND, the breakdown $B_i$ must be temporal; for ATTRIBUTION and DISTRIBUTION\_DIFFERENCE, $B_i$ must be categorical or ID-type; for OUTSTANDING\_VALUE, any breakdown type is valid. We report SVR overall and also by pattern type. This also ensures that breakdown-measure pairs are actionable for subgroup analysis.

\textit{Pattern Coverage} counts how many out of the four patterns appear with a valid structure. A system that only produces one or two patterns is not doing broad EDA.

The formula is:
\begin{equation}
\text{Pattern\_Coverage} = \frac{|\{p \in P : \exists i \in I, \text{type}(i) = p \land \text{valid}(i)\}|}{|P|}
\end{equation}

where $P = \{\text{TR, OV, AT, DD}\}$ are the four pattern types (TREND, OUTSTANDING\_VALUE, ATTRIBUTION, DISTRIBUTION\_DIFFERENCE), and $\text{valid}(i)$ means the insight passes SVR check.

\subsubsection{Group 3: Statistical Quality}

\textit{Statistical Significance} checks the percentage of insights that pass a statistical test at p < 0.05. We use pattern-specific tests: Mann-Kendall for trends \cite{mann}, z-test for outstanding values, Chi-square test for attribution \cite{mchugh}, and Kolmogorov-Smirnov test for distribution differences \cite{kolmogorov}. Effect sizes are measured using Kendall's $\tau$ for trends \cite{kendall}, z/(z+1) for outstanding values, Cramér's V for attribution \cite{cramer}, and KS statistic for distribution differences.

The formula is:
\begin{equation}
\text{Statistical\_Significance} = \frac{|\{i \in I : p\text{-value}(i) < 0.05\}|}{|I|}
\end{equation}

where $p\text{-value}(i)$ is computed using the pattern-specific test for insight $i$. Only structurally valid insights (SVR = 1) are included in this calculation.

\subsubsection{Group 4: Exploration Quality}

\textit{Subspace Rate} measures how many insights contain at least one subspace filter. A high value means the system does not only provide global insights but also explores conditional cases, such as ``within Region = West''.

The formula is:
\begin{equation}
\text{Subspace\_Rate} = \frac{|\{i \in I : |S_i| > 0\}|}{|I|}
\end{equation}

where $S_i$ is the set of subspace filters for insight $i$. Each filter is a (column, value) pair that restricts the data to a subset.

\textit{Score Uplift from Subspace} measures whether subspace insights have stronger patterns than global insights. We define two related quantities. The absolute mean score of subspace insights is:
\begin{equation}
x = \frac{1}{|I_S|}\sum_{i \in I_S} \text{score}(i)
\end{equation}
and the uplift delta relative to global insights is:
\begin{equation}
\Delta = \frac{1}{|I_S|}\sum_{i \in I_S} \text{score}(i) - \frac{1}{|I_G|}\sum_{i \in I_G} \text{score}(i)
\end{equation}
where $I_S = \{i \in I : |S_i| > 0\}$, $I_G = \{i \in I : |S_i| = 0\}$, and $\text{score}(i)$ is the pattern strength (Kendall's $\tau$ for TREND, effect size otherwise). We report $x$ in all tables since $|I_G|$ can be zero for some systems; $\Delta$ is reported per-dataset in supplementary results.

\textit{Simpson's Paradox Rate (SPR)} detects statistically significant pattern reversals between global and subspace insights \cite{simpson}. A reversal occurs when the pattern direction changes---for example, a positive trend globally becomes negative in a subspace. We use pattern-specific tests (Mann-Kendall for TREND, Chi-square contingency table for ATTRIBUTION, KS test for DISTRIBUTION\_DIFFERENCE) and require both global and subspace patterns to be statistically significant (p < 0.05).

The formula is:
\begin{equation}
\text{SPR} = \frac{|\{i \in I_S : \text{is\_paradox}(i) \land \text{is\_significant}(i)\}|}{|I_S|}
\end{equation}

where $\text{is\_paradox}(i)$ detects pattern reversal and $\text{is\_significant}(i)$ requires p < 0.05 for both global and subspace patterns. Note that SPR is computed over subspace insights only ($I_S$), so a system with fewer subspace insights has fewer opportunities to produce paradoxes; Subspace Rate should be read alongside SPR.

\subsubsection{Group 5: Alignment and Coherence}

\textit{Question-Insight Alignment} measures how close the question text is to the final insight text. We compute cosine similarity using sentence embeddings. For example, related sentence pairs on general text typically score 0.5 with MiniLM embeddings. This metric ensures that if QUIS and Baseline align similarly, differences elsewhere reflect insight quality rather than ISGEN favoring one system.

The formula is:
\begin{equation}
\text{Q-I\_Alignment} = \frac{1}{|I|} \sum_{i \in I} \text{cosine\_sim}(\text{emb}(Q_i), \text{emb}(T_i))
\end{equation}

where $Q_i$ is the question text, $T_i$ is the insight text, and $\text{emb}(\cdot)$ uses Sentence-BERT embeddings \cite{sbert}.

\textit{Reason-Insight Coherence} measures how close the reason field is to the final insight. This metric is only used for QUIS and the LLM Baseline because ONLYSTATS only uses template-based questions.

The formula is:
\begin{equation}
\text{R-I\_Coherence} = \frac{1}{|I|} \sum_{i \in I} \text{cosine\_sim}(\text{emb}(R_i), \text{emb}(T_i))
\end{equation}

where $R_i$ is the reason text for insight $i$. This metric is not applicable to ONLYSTATS.

\section{Experimental Setup}

We use three datasets summarized in Table~\ref{tab:datasets}: Adidas US Sales (Chaudhari, 2022) \cite{adidas}, IBM Employee Attrition (Subhash, 2017) \cite{ibm}, and Online Sales (Verma, 2024) \cite{online}. We use OpenAI's \texttt{gpt-5.4} for QUIS and the LLM Baseline. The original paper used Llama-3-70B-instruct; our model may not produce exact numbers but does not affect comparability since all systems share the same evaluation framework. We set \texttt{beam\_width = 100}, matching the original paper's configuration. For text similarity metrics, we use \texttt{all-MiniLM-L6-v2} from Sentence-BERT (Reimers and Gurevych, 2019) \cite{sbert}.

\section{Results}

\subsection{Overview}

{
\renewcommand{\arraystretch}{0.9}
\setlength{\belowcaptionskip}{-2pt}
\begin{table}[htbp]
\caption{Dataset Characteristics}
\begin{center}
\resizebox{\columnwidth}{!}{%
\begin{tabular}{|l|r|r|p{4cm}|}
\hline
\textbf{Dataset} & \textbf{Rows} & \textbf{Cols} & \textbf{Key Features} \\
\hline
Adidas US Sales & 9,648 & 13 & Temporal, categorical, numerical \\
Employee Attrition & 1,470 & 35 & Binary, ordinal, many categorical \\
Online Sales & 240 & 9 & Small, categorical, numerical \\
\hline
\end{tabular}}
\label{tab:datasets}
\end{center}
\end{table}
}

Table~\ref{tab:avg_results} summarizes average results across all three datasets. Detailed per-dataset results are shown in Tables~\ref{tab:adidas_results}--\ref{tab:online_results}.

{
\renewcommand{\arraystretch}{0.9}
\setlength{\belowcaptionskip}{-2pt}
\begin{table*}[htbp]
\caption{Average Evaluation Results Across Three Datasets}
\begin{center}
\begin{tabular}{|l|l|c|c|c|c|}
\hline
\textbf{Group} & \textbf{Metric} & \textbf{QUIS} & \textbf{Baseline} & \textbf{ONLYSTATS} & \textbf{Winner} \\
\hline
Correctness & 1. Faithfulness & \textbf{100.0\%} & \textbf{100.0\%} & \textbf{100.0\%} & Tie \\
Validity & 2. SVR (Overall) & \textbf{94.0\%} & 40.0\% & 90.8\% & \textbf{QUIS} \\
Coverage & 3. Pattern Coverage & 3--4/4 & 2--3/4 & 3--4/4 & ONLYSTATS \\
Statistics & 4. Statistical Significance & 46.4\% & 57.6\% & \textbf{58.0\%} & \textbf{ONLYSTATS} \\
Exploration (Qty) & 5. Subspace Rate & \textbf{84.4\%} & 37.4\% & 62.6\% & \textbf{QUIS} \\
Exploration (Qual) & 6. Score Uplift ($x$) & \textbf{1.067} & 0.974 & 0.818 & \textbf{QUIS} \\
Exploration (Qual) & 7. Simpson's Paradox Rate & 30.1\% & 18.9\% (0 sig) & \textbf{34.7\%} (0 sig) & \textbf{ONLYSTATS} \\
Control & 8. Q-I Alignment & 0.540 & \textbf{0.569} & N/A & \textbf{Baseline} \\
Control & 9. R-I Coherence & \textbf{0.526} & 0.514 & N/A & \textbf{QUIS} \\
\hline
\end{tabular}
\label{tab:avg_results}
\end{center}
\end{table*}
}

\subsection{Detailed Results by Dataset}

{
\renewcommand{\arraystretch}{0.9}
\setlength{\belowcaptionskip}{-10pt}
\begin{table}[!ht]
\caption{Results on Adidas US Sales (QUIS=99, Base=75, ONLY=113)}
\begin{center}
\resizebox{\columnwidth}{!}{%
\begin{tabular}{|l|c|c|c|}
\hline
\textbf{Metric} & \textbf{QUIS} & \textbf{Baseline} & \textbf{ONLYSTATS} \\
\hline
Faithfulness & \textbf{100\%} & \textbf{100\%} & \textbf{100\%} \\
SVR (Overall) & 99.0\% & 45.3\% & \textbf{96.5\%} \\
SVR - ATTRIBUTION & \textbf{100\%} & N/A & \textbf{100\%} \\
SVR - TREND & \textbf{100\%} & 48.5\% & \textbf{100\%} \\
Pattern Coverage & \textbf{4/4} & 3/4 & \textbf{4/4} \\
Statistical Significance & \textbf{83.4\%} & 73.2\% & 80.7\% \\
Subspace Rate & \textbf{86.9\%} & 42.7\% & 61.9\% \\
Score Uplift ($x$) & 0.885 & 0.796 & 0.715 \\
SPR & 27.9\% (0 sig) & 25.0\% (0 sig) & \textbf{42.9\%} (0 sig) \\
Q-I Alignment & \textbf{0.583} & 0.579 & N/A \\
R-I Coherence & \textbf{0.553} & 0.527 & N/A \\
\hline
\end{tabular}}
\label{tab:adidas_results}
\end{center}
\end{table}
}
\vspace{-4pt}

Table~\ref{tab:adidas_results} shows that QUIS leads on SVR (99.0\% vs 45.3\%) and Subspace Rate (86.9\%), while ONLYSTATS accomplishes the highest raw SPR (42.9\%). On Adidas all systems produce zero statistically significant paradoxes, proving that patterns on this dataset do not meet the strict dual-significance threshold.

{
\renewcommand{\arraystretch}{0.9}
\setlength{\belowcaptionskip}{-10pt}
\begin{table}[!ht]
\caption{Results on IBM Employee Attrition (QUIS=133, Base=81, ONLY=125)}
\begin{center}
\resizebox{\columnwidth}{!}{%
\begin{tabular}{|l|c|c|c|}
\hline
\textbf{Metric} & \textbf{QUIS} & \textbf{Baseline} & \textbf{ONLYSTATS} \\
\hline
Faithfulness & \textbf{100\%} & \textbf{100\%} & \textbf{100\%} \\
SVR (Overall) & 96.2\% & 37.0\% & \textbf{100.0\%} \\
Pattern Coverage & 3/4 & 3/4 & 3/4 \\
Statistical Significance & 20.0\% & \textbf{55.8\%} & 30.6\% \\
Subspace Rate & \textbf{87.2\%} & 33.3\% & 60.0\% \\
Score Uplift ($x$) & \textbf{1.574} & 1.079 & 0.506 \\
SPR & \textbf{30.2\%} (1 sig) & 0.0\% (0 sig) & 21.3\% (0 sig) \\
Q-I Alignment & 0.493 & \textbf{0.588} & N/A \\
R-I Coherence & 0.468 & \textbf{0.519} & N/A \\
\hline
\end{tabular}}
\label{tab:employee_results}
\end{center}
\end{table}
}
\vspace{-4pt}

Table~\ref{tab:employee_results} shows that QUIS achieves the strongest Score Uplift ($x$=1.574) and is the only system with significant paradoxes (SPR 30.2\%, 1 significant). ONLYSTATS achieves 21.3\% SPR with zero significant paradoxes, showing that exhaustive enumeration does not guarantee subspace quality on datasets with many weak global patterns.

{
\renewcommand{\arraystretch}{0.9}
\setlength{\belowcaptionskip}{-10pt}
\begin{table}[!ht]
\caption{Results on Online Sales (QUIS=106, Base=61, ONLY=114)}
\begin{center}
\resizebox{\columnwidth}{!}{%
\begin{tabular}{|l|c|c|c|}
\hline
\textbf{Metric} & \textbf{QUIS} & \textbf{Baseline} & \textbf{ONLYSTATS} \\
\hline
Faithfulness & \textbf{100\%} & \textbf{100\%} & \textbf{100\%} \\
SVR (Overall) & 86.8\% & 37.7\% & 77.2\% \\
Pattern Coverage & 3/4 & 2/4 & \textbf{4/4} \\
Statistical Significance & 35.9\% & 43.8\% & \textbf{62.8\%} \\
Subspace Rate & 79.2\% & 36.1\% & 65.8\% \\
Score Uplift ($x$) & 0.742 & \textbf{1.048} & 1.232 \\
SPR & 32.1\% (1 sig) & 31.8\% (0 sig) & \textbf{40.0\%} (0 sig) \\
Q-I Alignment & \textbf{0.543} & 0.539 & N/A \\
R-I Coherence & \textbf{0.557} & 0.497 & N/A \\
\hline
\end{tabular}}
\label{tab:online_results}
\end{center}
\end{table}
}
\vspace{-4pt}

Table~\ref{tab:online_results} shows that ONLYSTATS achieves the highest Statistical Significance (62.8\%) and Score Uplift ($x$=1.232), while the Baseline achieves the highest raw Score Uplift ($x$=1.048) due to its smaller insight set concentrating on globally stronger patterns. QUIS leads on Subspace Rate (79.2\%). All three systems achieve raw SPR (31--40\%); QUIS produces one significant paradox, while ONLYSTATS and Baseline find none.

\subsection{Cross-Dataset Analysis}

Our metrics are interconnected: structural validity enables subspace exploration, which affects subspace quality. This allows us to trace how design choices propagate through the pipeline

\subsubsection{Faithfulness}

All three systems achieve 100\% faithfulness across all datasets, confirming that reported values are verifiable and that QUGEN does not sacrifice correctness for richer question generation.

\subsubsection{Structural Validity Rate}

SVR shows the largest and most consistent difference across datasets. QUIS achieves 94.0\% average SVR, while the LLM Baseline achieves only 40.0\% (gap of 54 percentage points). ONLYSTATS achieves 91.2\% by constraining its statistics-based selection to valid column types, ensuring breakdown columns match pattern requirements.

The Baseline's core problem is systematic column-type mismatch: it achieves 0\% ATTRIBUTION SVR on Adidas and Online Sales, always selecting numeric columns (e.g., \texttt{Total Sales}, \texttt{MonthlyIncome}) as the breakdown instead of categorical ones. For TREND, it uses categorical breakdowns in 51.5\% of Adidas cases (e.g., \texttt{Region}, \texttt{Product}) instead of temporal columns. QUIS avoids this because QUGEN produces structured cards with explicit, programmatically validated breakdown fields. Notably, the Baseline's prompt explicitly instructs the LLM to use categorical or temporal columns as breakdowns for TREND, ATTRIBUTION, and DISTRIBUTION\_DIFFERENCE patterns, and to avoid numerical columns. However, the LLM fails to follow these instructions consistently, resulting in the observed structural validity issues. This highlights a key limitation of relying solely on prompt-based constraints for enforcing structural correctness in automated EDA systems.

\subsubsection{Pattern Coverage}

All systems miss TREND on Employee Attrition and Online Sales due to limited temporal columns. The Baseline misses ATTRIBUTION on Adidas and Online Sales because its numeric breakdowns are structurally invalid. This limitation stems from dataset characteristics: lack of temporal columns prevents TREND analysis, while the Baseline's structural validity problems prevent ATTRIBUTION generation where categorical breakdowns are required.

\subsubsection{Subspace Exploration}

Subspace Rate is where QUIS shows the most advantage. QUIS achieves 84.4\% average subspace rate (range: 79.2\%--87.2\%), compared with 37.4\% for the Baseline (range: 33.3\%--42.7\%) and 62.6\% for ONLYSTATS (range: 60.0\%--65.8\%).

Since QUIS and ONLYSTATS share the same ISGEN module, the 21.8-point average gap between them isolates the contribution of QUGEN's insight cards over exhaustive enumeration. The 47.0-point gap versus the Baseline reflects its structural problems: invalid breakdowns cannot be explored in subspaces. ONLYSTATS reaches its highest subspace rate on Online Sales (65.8\%) because exhaustive enumeration produces many valid (B,\,M) pairs of this small 9-column dataset, giving ISGEN more entry points.

Score Uplift is positive on average for all systems (QUIS: 1.067, Baseline: 0.974, ONLYSTATS: 0.818), measured as the mean pattern strength of subspace insights ($x$). QUIS shows the strongest uplift on Employee Attrition ($x$=1.574) and leads on average across all datasets. On Online Sales, ONLYSTATS achieves the highest Score Uplift ($x$=1.232), while the Baseline achieves $x$=1.048 due to its smaller insight set concentrating on globally stronger patterns. SPR follows a similar trend: ONLYSTATS achieves the highest average (34.7\%) but with moderate variance across datasets (21.3\%--42.9\%), while QUIS achieves 30.1\% with more consistency (27.9\%--32.1\%). The Baseline finds no significant paradoxes across any dataset, confirming that structurally invalid insights cannot produce meaningful subspace patterns.

\subsubsection{Statistical Significance}

Statistical significance is dataset-dependent (QUIS: 46.4\%, Baseline: 57.6\%, ONLYSTATS: 58.0\% on average). QUIS leads on Adidas (83.4\%) but is less efficient on Employee Attrition (20.0\%), where binary/categorical variables make simple aggregations easier to pass significance tests. This dependency reflects inherent difficulty: temporal datasets with continuous measures favor QUIS's focused subspace search, while categorical-heavy datasets benefit from exhaustive enumeration.

\subsubsection{Control Metrics}

Q-I Alignment is similar across QUIS (0.540) and Baseline (0.569), confirming that ISGEN answers both systems comparably---differences come from card quality, not retrieval. R-I Coherence slightly favors QUIS on average (0.526 vs 0.514), though the Baseline scores higher on Employee Attrition (0.519 vs 0.468). This reversal is consistent with the Baseline's higher statistical significance on that dataset (55.8\% vs 20.0\%): when Baseline insights happen to be statistically stronger, their text tends to be more specific and that are better aligned with the reason. Both metrics are not reported for ONLYSTATS (template-based questions only).

\section{Discussion}

\subsection{What the QUIS Pipeline Adds}

QUGEN's primary contribution is structural correctness: its schema-aware cards with explicit breakdown fields eliminate the column-type mismatches that plague the free-form LLM Baseline (0\% ATTRIBUTION SVR on two of three datasets, 40.0\% overall). ONLYSTATS also achieves high SVR by constraining enumeration to valid column types, showing that the constraint itself matters, not the generation method.

For subspace exploration, QUGEN's contribution is indirect but measurable: semantically grounded cards make ISGEN more focused entry points, resulting in higher Subspace Rate and more consistent SPR across datasets. ONLYSTATS, though using the same ISGEN module, shows greater variance in SPR (21.3\%--42.9\%), suggesting that exhaustive enumeration produces many (B, M) pairs where global patterns are too weak for subspace search to improve. The LLM Baseline's zero significant paradoxes are mainly associated with its low structural validity, consistent with the view that valid breakdown columns are necessary for meaningful subspace patterns.

\subsection{Limitations}

There are some limitations in our study. First, we use OpenAI's \texttt{gpt-5.4} instead of the original Llama-3-70B-instruct, so the exact numbers may be different from the original paper. Second, our embedding-based metrics depend on \texttt{all-MiniLM-L6-v2}. This model is useful, but it may not capture all meanings in short technical texts. Third, while we test on three datasets with varying characteristics, broader coverage would further validate our findings. Fourth, with only three datasets, cross-system significance tests (e.g., paired Wilcoxon) are not feasible, so metric differences should be interpreted as descriptive rather than statistically confirmed. Nevertheless, the observed gaps are large and consistent across datasets. For example, QUIS shows a 54-point SVR improvement (94.0\% vs 40.0\%) and a 47-point Subspace Rate improvement (84.4\% vs 37.4\%) over the Baseline. These differences are consistent with the structural behavior of the systems: invalid breakdown columns limit subspace exploration, which explains the Baseline's low Subspace Rate and absence of significant paradoxes. Overall, the consistency of these patterns suggests systematic design differences rather than random variation.

\section{Conclusion}

We replicated QUIS and built an automatic, human-free evaluation framework with 9 metrics covering correctness, structural validity, statistical quality, and subspace exploration. Across three datasets, results show three outstanding findings. First, the LLM Baseline has a systematic structural validity problem (37.0\%--45.3\% SVR) caused by column-type mismatches, while QUIS achieves 86.8\%--99.0\% across datasets. Second, QUGEN provides higher quality input cards that enable ISGEN to discover more subspace insights (84.4\% Subspace Rate vs 37.4\%) with stronger and more consistent patterns (SPR 30.1\%, Score Uplift $x$=1.067). Third, statistical significance is dataset-dependent and should not be interpreted in isolation. These findings have two practical implications: LLM-based EDA systems should validate breakdown column types before output, and insight search engines benefit from semantically grounded input cards over exhaustive enumeration alone. Our main contribution is the evaluation framework itself, which is a reusable tool for comparing any AutoEDA system in a consistent and reproducible way.

\begin{thebibliography}{00}
\bibitem{quis} A. Manatkar, A. Akella, P. Gupta, and K. Narayanam, ``QUIS: Question-guided Insights Generation for Automated Exploratory Data Analysis,'' arXiv:2410.10270, 2024.
\bibitem{quickinsights} R. Ding, S. Han, Y. Xu, H. Zhang, and D. Zhang, ``QuickInsights: Quick and Automatic Discovery of Insights from Multi-Dimensional Data,'' in \textit{SIGMOD}, pp. 317--332, 2019.
\bibitem{metainsight} P. Ma, R. Ding, S. Han, and D. Zhang, ``MetaInsight: Automatic Discovery of Structured Knowledge for Exploratory Data Analysis,'' in \textit{SIGMOD}, pp. 1262--1274, 2021.
\bibitem{insightpilot} P. Ma, R. Ding, S. Wang, S. Han, and D. Zhang, ``InsightPilot: An LLM-Empowered Automated Data Exploration System,'' in \textit{EMNLP System Demonstrations}, pp. 346--352, 2023.
\bibitem{seedb} M. Vartak, S. Rahman, S. Madden, A. Parameswaran, and N. Polyzotis, ``SEEDB: Efficient Data-Driven Visualization Recommendations to Support Visual Analytics,'' in \textit{VLDB}, pp. 2182--2193, 2015.
\bibitem{zenvisage} T. Siddiqui, A. Kim, J. Lee, K. Karahalios, and A. Parameswaran, ``Effortless Data Exploration with zenvisage: An Expressive and Interactive Visual Analytics System,'' in \textit{VLDB}, pp. 457--468, 2016.
\bibitem{voyager} K. Wongsuphasawat, D. Moritz, A. Anand, J. Mackinlay, B. Howe, and J. Heer, ``Voyager: Exploratory Analysis via Faceted Browsing of Visualization Recommendations,'' \textit{IEEE Transactions on Visualization and Computer Graphics}, vol. 22, no. 1, pp. 649--658, 2016.
\bibitem{voder} A. Srinivasan and J. Stasko, ``Augmenting Visualizations with Interactive Data Facts,'' \textit{IEEE Transactions on Visualization and Computer Graphics}, vol. 25, no. 1, pp. 672--681, 2018.
\bibitem{bach} J. Bach, ``Subgroup Discovery with Small and Alternative Feature Sets,'' \textit{Proc. ACM Manag. Data}, vol. 3, no. 3, Article 221, 2025.
\bibitem{kruskal} W. H. Kruskal and W. A. Wallis, ``Use of ranks in one-criterion variance analysis,'' \textit{Journal of the American Statistical Association}, vol. 47, no. 260, pp. 583--621, 1952.
\bibitem{simpson} E. H. Simpson, ``The Interpretation of Interaction in Contingency Tables,'' \textit{Journal of the Royal Statistical Society, Series B}, vol. 13, no. 2, pp. 238--241, 1951.
\bibitem{teng} X. Teng and Y.-R. Lin, ``De-paradox Tree: Breaking Down Simpson's Paradox via A Kernel-Based Partition Algorithm,'' arXiv:2603.02174, 2026.
\bibitem{mann} H. B. Mann, ``Nonparametric Tests Against Trend,'' \textit{Econometrica}, vol. 13, pp. 245--259, 1945.
\bibitem{kendall} M. G. Kendall, \textit{Rank Correlation Methods}, 4th ed. Griffin, 1975.
\bibitem{cramer} H. Cramér, \textit{Mathematical Methods of Statistics}. Princeton University Press, 1946.
\bibitem{kolmogorov} A. Kolmogorov, ``Sulla determinazione empirica di una legge di distribuzione,'' \textit{Giornale dell'Istituto Italiano degli Attuari}, vol. 4, pp. 83--91, 1933.
\bibitem{mchugh} M. L. McHugh, ``The Chi-square Test of Independence,'' \textit{Biochemia Medica}, vol. 23, no. 2, pp. 143--149, 2013.
\bibitem{sbert} N. Reimers and I. Gurevych, ``Sentence-BERT: Sentence Embeddings using Siamese BERT-Networks,'' in \textit{EMNLP}, pp. 3982--3992, 2019.
\bibitem{adidas} H. Chaudhari, ``Adidas Sales Dataset,'' Kaggle, 2022.
\bibitem{ibm} P. Subhash, ``IBM HR Analytics Employee Attrition \& Performance,'' Kaggle, 2017.
\bibitem{online} S. Verma, ``Online Sales Dataset - Popular Marketplace Data,'' Kaggle, 2024.
\end{thebibliography}

\end{document}
